% 张三慧第4版第二卷第22章《光的干涉》原书思考题22.1--22.12
% 教材V2 PDF物理页261--262（印刷244--245）；独立于正式习题22.1--22.23。
\section*{第22章\quad 光的干涉：思考题}
\paragraph{思考题 22.1}
用白光源经双缝实验，第一缝后放红滤光片、
第二缝后放绿滤光片，还能观察稳定的双缝干涉条纹吗？
解释光谱重叠的影响。
\paragraph{思考题 22.2}
对比图22.30中钠光灯和激光器照明双缝
两种方案，是否一定都可观察清晰干涉条纹？
从光源空间相干性与时间相干性分别判断。
\paragraph{思考题 22.3}
为什么两点相干源形成的水波干涉图案中，
相同相位差的条纹常沿双曲线分布？
\paragraph{思考题 22.4}
把两个顶角很小的棱镜底边拼成
Fresnel双棱镜，如何从一个原始光源生成
两个相干的虚光源？其两光束的重叠区在哪里？
\paragraph{思考题 22.5}
两束光的强度分别为 \(I_1,I_2\)，
若彼此相干，重叠时总强度怎样计算？
若两束光不相干，时间平均总强度怎样计算？
\paragraph{思考题 22.6}
Young双缝实验中：
(1) 增大缝距 \(d\) 后条纹间距如何变化？
(2) 横向移动有限距离的原光源，
条纹中心如何变化？
(3) 把原光源的发光区域逐渐加宽，
条纹可见度有什么变化？
\paragraph{思考题 22.7}
利用一束光的时间延迟两路产生干涉
和利用同一波前不同位置产生干涉，
分别测量了光的哪种相干性？
给出相干时间/相干长度和空间相干宽度的物理含义。
\paragraph{思考题 22.8}
肥皂泡表面的彩色来自什么干涉？
为什么泡壁足够薄时可能出现黑色区域，
并在破裂前消失彩色？
\paragraph{思考题 22.9}
两块玻璃板构成空气劈尖（图22.32），
(1) 使上表面缓慢整体平移、厚度处处增加时
等厚干涉条纹如何移动？
(2) 增大劈尖角时条纹间距如何变化？
\paragraph{思考题 22.10}
用普通单色光照射一块两表面夹角很小、
但总体厚度很大的玻璃劈尖，
为什么可能观察不到干涉条纹？
\paragraph{思考题 22.11}
利用相距数百至数千千米、同步观测的射电望远镜
如何通过相位差、时间延迟测量
地球自转参数及大陆板块的相对移动？
\paragraph{思考题 22.12}
Young双缝上方缝后插入一块折射率 \(n>1\)
的薄云母片，厚度为 \(t\)。
双缝干涉条纹的间距是否变化？
原中央零级明纹向哪一侧移动，移动多少？
